\section{Computational Fairness}
\label{sec:main-idea}

Our notion of computational fairness starts with two players, Alice and Bob. Alice runs an algorithm $\mathcal{A}$ that makes decisions based on some input.  For example, Alice may be an employer using $\mathcal{A}$ to decide who to hire. Specifically, $\mathcal{A}$ takes a data set $D$ with protected attribute $X$ and unprotected attributes $Y$ and makes a (binary) decision $C$. By law, Alice is not allowed to use $X$ in making decisions, and claims to use only $Y$. It is Bob's job to verify that on the data $D$, Alice's algorithm $\mathcal{A}$ is not liable for a claim of disparate impact. 

\vspace*{0.1in}\mypara{Trust model.} We assume that Bob does not have access to
$\mathcal{A}$. Further, we assume that Alice has good intentions: specifically,
that Alice is not secretly using $X$ in $\mathcal{A}$ while lying about
it. While assuming Alice is lying about the use of $X$ might be more plausible, it is much harder to detect. More importantly, from a functional perspective, it does not matter whether Alice uses $X$ explicitly or uses proxy attributes $Y$ that have the same effect: this is the core message from the Griggs case that introduced the doctrine of disparate impact. In other words, our \emph{certification} process is indifferent to Alice's intentions, but our \emph{repair} process will assume good faith.

We summarize our main idea with the following intuition:
\begin{quote}
\emph{  If Bob cannot predict $X$ given the other attributes of $D$, then $\mathcal{A}$ is fair with respect to Bob on $D$. % More generally, if the balanced error rate of any predictor is at least $\epsilon$, then the algorithm is $\epsilon$-fair. 
}\end{quote}

% Bob can be thought of as an ensemble of all possible models, where the bias of $D$ on the most biased model is determined to be the true bias.  (This also hints at an implementation of Bob in which Bob runs a large collection of models and measures the bias in the results.)

\subsection{Predictability and Disparate Impact}

We now present a formal definition of predictability and link it to the legal notion of disparate impact.  Recall that $D = (X,Y,C)$ where $X$ is the protected attribute, $Y$ is the remaining attributes, and $C$ is the class outcome to be predicted.

The basis for our formulation is a procedure that predicts $X$ from $Y$. We would like a way to measure the quality of this predictor in a way that a) can be optimized using standard predictors in machine learning and b) can be related to $\lrplus$. The standard notions of accuracy of a classifier fail to do the second (as discussed earlier) and using $\lrplus$ directly fails to satisfy the first constraint. 

The error measure we seek turns out to be the \emph{balanced error rate} \ber. 

\begin{defn}[BER]
  Let $f : Y \to X$ be a predictor of $X$ from $Y$. The \emph{balanced error rate} \ber of $f$ on distribution $\mathcal{D}$ over the pair $(X, Y)$ is defined as the (unweighted) average class-conditioned error of $f$. In other words, 
\[ \ber(f(Y), X) = \frac{\Pr[f(Y) = 0 | X = 1] + \Pr[f(Y) = 1 | X = 0] }{2} \]
\end{defn}

\begin{defn}[{\small Predictability}]
$X$ is said to be \emph{$\epsilon$-predictable} from $Y$ if there exists a function $f:Y \rightarrow X$ such that
\[\ber(f(Y), X) \le\epsilon. \]
\end{defn}

This motivates our definition of $\epsilon$-fairness, as a data set that is \emph{not} predictable.
\begin{defn}[$\epsilon$-fairness]
A data set $D = (X, Y, C)$ is said to be \emph{$\epsilon$-fair} if for any
classification algorithm $f \colon Y \to X$ 
\[\ber(f(Y), X) > \epsilon \]
with (empirical) probabilities estimated from $D$.
\end{defn}

Recall the definition of disparate impact from Section \ref{sec:note-meas-disp}.  We will be interested in examining the potential disparate impact of a classifier $g \colon Y \rightarrow C$ and will consider the value $\di(g) = 1 / \lrplus(g(Y), X)$ as it relates to the threshold $\tau$.  Where $g$ is clear from context, we will refer to this as $\di$.

The justification of our definition of fairness comes from the following theorem:
\begin{theorem}
\label{sec:conn-pred-disp-1}
A data set is $(1/2-\beta/8)$-predictable if and only if it admits disparate impact, where $\beta$ is the fraction of elements in the minority class $(X = 0)$ that are selected $(C = 1)$. 
\end{theorem}

\begin{proof}
We will start with the direction showing that disparate impact implies predictability.  Suppose that there exists some function $g:Y \rightarrow C$ such that $\lrplus(g(y), c) \geq \frac{1}{\tau}$.  We will create a function $\psi:C \rightarrow X$ such that $\ber(\psi(g(y)), x) < \epsilon$ for $(x,y) \in D$. Thus the combined predictor $f = \psi \circ g$ satisfies the definition of predictability. 

Consider the confusion matrix associated with $g$, depicted in \autoref{tab:di-confusion}.
\begin{table}[htbp]
  \centering
  \begin{tabular}{r|c|c}
   Prediction & $X = 0$ & $X = 1$ \\ \hline
   $g(y) = \no  $ & a       & b       \\
   $g(y) = \yes $ & c       & d
  \end{tabular}
  \caption{Confusion matrix for $g$\label{tab:di-confusion}}
\end{table}
Set $\alpha \triangleq
\frac{b}{b+d}$ and $\beta\triangleq \frac{c}{a+c}$. Then we can write 
$ \lrplus(g(y), X) = \frac{1-\alpha}{\beta}$ and $\di(g) = \frac{\beta}{1 - \alpha}$. 

We define the \emph{purely biased} mapping $\psi \colon C \to X$ as $\psi(\yes) = 1$ and $\psi(\no) = 0$. Finally, let $\phi \colon Y \to X = \psi \circ g$. The confusion matrix for $\phi$ is depicted in \autoref{tab:ber-confusion}. Note that the confusion matrix for $\phi$ is identical to the matrix for $g$. 
\begin{table}[htbp]
  \centering
  \begin{tabular}{r|c|c}
   Prediction & $X = 0$ & $X = 1$ \\ \hline
   $\phi(Y) = 0  $ & a       & b       \\
   $\phi(Y) = 1  $ & c       & d
  \end{tabular}
  \caption{Confusion matrix for $\phi$\label{tab:ber-confusion}}
\end{table}

We can now express $\ber(\phi)$ in terms of this matrix. Specifically, 
$ \ber(\phi) = \frac{\alpha + \beta}{2}$.

\vspace*{0.1in}
\mypara{Representations.} 
We can now express contours of the $\di$ and $\ber$ functions as curves in the unit square $[0,1]^2$. 
Reparametrizing $\pi_1 = 1 - \alpha$ and $\pi_0 = \beta$, we can express the error
measures as $  \di(g) =  \frac{\pi_0}{\pi_1} $ and $ \ber(\phi) = \frac{1 + \pi_0 - \pi_1}{2} $


As a consequence, any classifier $g$ with $\di(g) = \delta$ can be represented
in the $[0,1]^2$ unit square as the line $\pi_1 = \pi_0/\delta$. Any classifier $\phi$
with $\ber(\phi) = \epsilon$ can be written as the function $\pi_1 = \pi_0 + 1 - 2\epsilon$. 


Let us now fix the desired \di\ threshold $\tau$, corresponding to the line $\pi_1
=\pi_0/\tau$. Notice that the region $\{(\pi_0,\pi_1)\mid \pi_1 \ge \pi_0/\tau\}$ is the region
where one would make a finding of disparate impact (for $\tau = 0.8$).

Now given a classification that admits a finding of disparate impact, we can
compute $\beta$. Consider the point $(\beta, \beta/\tau)$ at which the line $\pi_0 = \beta$ intersects the
\di curve $\pi_1 = \pi_0/\tau$. This point lies on the \ber contour $(1 +
\beta-\beta/\tau)/2 = \epsilon$, yielding 
$ \epsilon = 1/2 - \beta(\frac{1}{\tau}-1)/2 $
In particular, for the \di threshold of $\tau = 0.8$, the desired \ber threshold is
\[\epsilon = \frac{1}{2} - \frac{\beta}{8} \]
and so disparate impact implies predictability.

With this infrastructure in place, the other direction of the proof is now easy.  To show that predictability implies disparate impact, we will use the same idea of a purely biased classifier.  Suppose there is a function $f \colon Y \to X$ such that $\ber(f(y), x) \leq \epsilon$.  Let $\psi^{-1} \colon X \to C$ be the inverse purely biased mapping, i.e. $\psi^{-1}(1) = \yes$ and $\psi^{-1}(0) = \no$.  Let $g \colon Y \to C = \psi^{-1} \circ f$.  Using the same representation as before, this gives us
$ \pi_1 \geq 1 + \pi_0 - 2  \epsilon$ and therefore
\[ \frac{\pi_0}{\pi_1} \leq \frac{\pi_0}{1 + \pi_0 - 2  \epsilon} = 1 - \frac{1 - 2 \epsilon}{\pi_0 + 1 - 2 \epsilon} \]
Recalling that $\di(g) = \frac{\pi_0}{\pi_1}$ and that $\pi_0 = \beta$ yields
$ \di(g) \leq 1 - \frac{1 - 2 \epsilon}{\beta + 1 - 2\epsilon} = \tau$.
For $\tau = 0.8$, this again gives us a desired \ber threshold of
$\epsilon = \frac{1}{2} - \frac{\beta}{8}$.
\end{proof}

Note that as $\epsilon$ approaches $1/2$ the bound tends towards
the trivial (since any binary classifier has \ber at most $1/2$). In
other words, as $\beta$ tends to $0$, the bound becomes vacuous. 

This points to an interesting line of attack to evade a disparate impact finding. Note that
$\beta$ is the (class conditioned) rate at which members of the protected class
are selected. Consider now a scenario where a company is being investigated for
discriminatory hiring practices. One way in which the company might defeat such
a finding is by interviewing (but not hiring) a large proportion of applicants
from the protected class. This effectively drives $\beta$ down, and the
observation above says that in this setting their discriminatory practices will
be harder to detect, because our result can not guarantee that a classifier will
have error significantly less than $0.5$. 

Observe that in this analysis we use an extremely weak classifier to prove the
existence of a relation between predictability and disparate impact. It is
likely that using a better classifier (for example the Bayes optimal classifier or even a classifier that optimizes \ber) might yield a stronger relationship between the two notions. 

\vspace*{0.1in}\mypara{Dealing with uncertainty.} In general, $\beta$ might be
hard to estimate from a fixed data set, and in practice we might only know that
the true value of $\beta$ lies in a range $[\beta_\ell, \beta_u]$. Since the
\ber threshold varies monotonically with $\beta$, we can merely use $\beta_\ell$
to obtain a conservative estimate. 

Another source of uncertainty is in the \ber estimate itself. Suppose that our
classifier yields an error that lies in a range $[\gamma, \gamma']$. Again,
because of monotonicity, we will obtain an interval of values $[\tau, \tau']$
for \di. Note that if (using a Bayesian approach) we are able to build a
distribution over \ber, this distribution will then transfer over to the \di
estimate as well. 

\subsection{Certifying (lack of) DI with SVMs}
\label{sec:certifying-lack-of}

The above argument gives us a way to determine whether a data set is potentially
amenable to disparate impact (in other words, whether there is insufficient
information to detect a protected attribute from the provided data).

\looseness-1 \vspace*{0.1in}\mypara{Algorithm.}
We run a classifier that optimizes \ber on the given data set, attempting to
predict the protected attributes $X$ from the remaining attributes $Y$.  Suppose
the error in this prediction is $\epsilon$. Then using the estimate of $\beta$
from the data, we can substitute this into the equation above and obtain a
threshold $\epsilon'$.  If $\epsilon' < \epsilon$, then we can declare the data
set free from disparate impact. 

\looseness-1 Assume that we have an optimal classifier with respect to \ber. Then we know
that all classifiers will incur a \ber of at least $\epsilon$. By Theorem
\ref{sec:conn-pred-disp-1}, this implies that no classifier on $D$ will exhibit
disparate impact, and so our certification is correct. 

The only remaining question is what classifier is used by this algorithm.  The usual way to incorporate class sensitivity into a classifier is to use different costs for misclassifying points in different classes. A number of class-sensitive cost measures fall into this framework, and there are algorithms for optimizing these measures (see \cite{menon2013statistical} for a review), as well as a general (but expensive) method due to Joachims that does a clever grid search over a standard SVM to optimize a large family of class-sensitive measures\cite{joachims2005support}. Oddly, \ber is not usually included among the measures studied. 

Formally, as pointed out by Zhao et al\cite{zhao2013beyond}, \ber is not a
cost-sensitive classification error measure because the weights assigned to
class-specific misclassification depend on the relative class sizes (so they can
be normalized). However, for any given data set we know the class sizes and can
reweight accordingly. We adapt a standard hinge-loss SVM to incorporate
class-sensitivity and optimize for (regularized) \ber. This adaptation is
standard, and yields a cost function that can be optimized using AdaBoost. %For
%more details, see the extended version of this paper \cite{arxiv-version}.

  We illustrate this by showing how to adapt a standard hinge-loss SVM to instead optimize \ber\footnote{This approach in general is well known; we provide a detailed derivation here for clarity.}. Consider the standard soft-margin (linear) SVM:
 \begin{equation}
 \min_{\vec{w},\boldsymbol{\xi},b}\frac{1}{2}\|\vec{w}\|^2 + \frac{C}{n}\sum_{j=1}^n\xi_j
 \end{equation}
such that
 \begin{equation} 
 y_j(\innerprod{\vec{w}}{\vec{x}_j} + b)\geq 1-\xi_j\quad\text{and}\quad \xi_j\geq 0\label{eq:1}.
 \end{equation}
 The constraints in \eqref{eq:1} are equivalent to the following:
 \begin{equation}
 \xi_j \geq \max\{1 - y_j(\innerprod{\vec{w}}{\vec{x}_j} + b), 0\},
 \end{equation}
 in which the right side is the hinge loss.

 Minimizing the \emph{balanced error rate} (BER) would result in the following optimization:
 \begin{equation}
   \min_{\vec{w},b}\frac{1}{2}\|\vec{w}\|^2 + 
   C\left( 
     \frac{1}{2n^+}\sum_{\{j|y_j=1\}}(1-\hat{y}_j) + 
     \frac{1}{2n^-}\sum_{\{j|y_j=-1\}}(1+\hat{y}_j) 
   \right),
 \end{equation}
 where $n^+ = |\{j\mid y_j=1\}|$ and $n^- = |\{j\mid y_j=-1\}|$ (we'll use $n_j$ to refer to the correct count for the corresponding $y_j$). 
 This optimization is NP-hard for the same reasons as minimizing the 0-1 loss is, so we will relax it to hinge loss as usual:
 \begin{align}
   &\min_{\vec{w},b}\ \frac{1}{2}\|\vec{w}\|^2 + 
   C\left( 
     \frac{1}{2n^+}\sum_{\{j|y_j=1\}}\xi_j + 
     \frac{1}{2n^-}\sum_{\{j|y_j=-1\}}\xi_j 
   \right),
 \\=&\min_{\vec{w},b}\ \frac{1}{2}\|\vec{w}\|^2 + 
   \frac{C}{n} \sum_j\frac{n}{2n_j}\xi_j
 \\=&\min_{\vec{w},b}\ \frac{1}{2}\|\vec{w}\|^2 + 
   \frac{C}{n} \sum_j D_j\xi_j,\label{eq:2}
 \end{align}
 where $\xi_j$ has the same constraints as in \eqref{eq:1}. 
 \eqref{eq:2} now has the same form as an SVM in AdaBoost, so we can use existing techniqes to solve this SVM. 
 Note also that $\sum_j D_j = 1$, making the $D_j$ a distribution over $[1..n]$ -- this means that we can use an AdaBoost formulation without needing to adjust constants.




% For $\gamma > 0$, let the \emph{approximate} \di\ threshold be $\tau' = (1+\gamma)\tau$, corresponding to the line $y = x/\tau'$. 

% We now invoke \autoref{cone-cover} with $u(x) = x/\tau$, $\ell(x) = x/\tau'$, and an arbitrary $x_0$. Let the resulting \pshiftID $h_k(x)$ consist of the partitioned range $[x_0, x_1, \ldots, x_k]$ and parameters $c_0, c_1, \ldots, c_{k-1}$. 





%%% Local Variables:
%%% mode: latex
%%% TeX-master: "bias"
%%% End:
